% Generated draft: requires mathematical and source review.
\documentclass[11pt]{article}
\usepackage[margin=25mm]{geometry}
\usepackage{fontspec}
\setmainfont{Latin Modern Roman}
\usepackage{amsmath,amssymb,mathtools,unicode-math}
\setmathfont{Latin Modern Math}
\usepackage{enumitem,needspace,xurl,hyperref}
\setlength{\parindent}{0pt}
\setlength{\parskip}{5pt}
\setlength{\emergencystretch}{3em}
\newcommand{\N}{\mathbb{N}}
\newcommand{\Z}{\mathbb{Z}}
\newcommand{\Q}{\mathbb{Q}}
\newcommand{\R}{\mathbb{R}}
\newcommand{\C}{\mathbb{C}}
\newcommand{\F}{\mathbb{F}}
\newcommand{\A}{\mathbb{A}}
\newcommand{\PP}{\mathbb{P}}
\newcommand{\E}{\mathbb{E}}
\newcommand{\eps}{\varepsilon}
\newcommand{\dd}{\,\mathrm{d}}
\begin{document}
\section*{UnsolvedMath ID: 30000810}
\textbf{Source problem code:} catalogue entry
\subsection{Definitions and mathematical statement}
A Moufang set consists of a set $X$ with at least three elements and subgroups $U_x\leq\operatorname{Sym}(X)$, one for each $x\in X$, such that $U_x$ fixes $x$ and acts sharply transitively on $X\setminus\{x\}$. The group generated by the root groups permutes them by conjugation according to its action on their fixed points.

In the usual coordinates $M(U,\tau)$, write the root group $U$ additively, without assuming it is abelian, and identify $X=U\cup\{\infty\}$. The root group at infinity acts by right translations $x\mapsto x+a$ on $U$ and fixes infinity. The permutation $\tau$ interchanges $0$ and $\infty$ and gives the associated opposite root group. Here $0$ is the identity of $U$ and $-a$ its group inverse. The special condition in these coordinates is
\[
\tau(-a)=-\tau(a)\qquad(a\in U\setminus\{0\}).
\]
The source asks whether every special Moufang set has an abelian root group:
\[
M(U,\tau)\text{ special}\quad\Longrightarrow\quad
\forall a,b\in U,\quad a+b=b+a.
\]
The additive notation in the hypotheses is purely notation and does not build the desired commutativity into the definition. All root groups are conjugate, so the question for $U$ is the corresponding question for every root group. No finiteness, element-order, field or topology hypothesis is imposed.
\subsection{Short English statement}
Must the root groups of every special Moufang set be abelian? The special coordinate identity is stated without assuming commutativity of the additively written group.
\subsection{Sources}
\begin{itemize}
\item Supplied report: \emph{Permutation Groups} (2008), Oberwolfach Reports, DOI 10.4171/owr/2007/37, pp. 2173--2236; proposed year 2007. \url{https://doi.org/10.4171/owr/2007/37}.
\item Supplied reference labelled ``Segev partial theorem'': \url{https://www.math.bgu.ac.il/~yoavs/venice/ys0856.pdf}.
\item Supplied reference labelled ``abelian centralizers'': \url{https://arxiv.org/abs/1009.2572}.
\item Local UnsolvedMath registry, permanent ID 30000810, catalogue entry; created 25 August 2026, updated 16 September 2026. Its corrected\_verified note describes editorial cleanup. The supplied 21 August 2026 literature triage labels the abelian-root-group conjecture open, reporting substantial special cases, element-order restrictions and abelianity of centralizers rather than a general theorem. No PDF access or independent updated assessment is asserted here.
\end{itemize}
\end{document}
